\NSPage{043}%
smoothly in \NSi{NS-F-P043-001}{\eta}. Their inverse bounds are finite because
\NSi{NS-F-P043-002}{\lambda}, the radii, and all bump functions were fixed before
\NSi{NS-F-P043-003}{N}.

Let \NSi{NS-F-P043-004}{\nu_k} be the norm of multiplication by
\NSi{NS-F-P043-005}{B^{-1}} on
\NSi{NS-F-P043-006}{C^k([-1,1])}, and let \NSi{NS-F-P043-007}{q_k} be the
bilinear norm of
\NSi{NS-F-P043-008}{\mathcal Q:C^k\times C^k\to C^k}, as in Lemma~\ref{lem:4.7}.
Since \NSi{NS-F-P043-009}{\lVert d_N\rVert_{C^k}\le D_k/N}, it suffices to
choose
\NSFormula{NS-F-P043-010}{display}{none}
\[
N\ge 8\max_{k=0,2}\nu_k^2q_kD_k
\]
to obtain a smooth solution of \eqref{eq:4.42} with
\NSi{NS-F-P043-011}{\lVert c\rVert_{C^2}\le 2\nu_2D_2/N}. The fixed bump
shapes then give
\NSFormula{NS-F-P043-012}{display}{none}
\[
\max_{r\le 1}\lVert \partial_X^r(u_c,e_c)\rVert_2\le C/N.
\]
Integrating the five correction integrands in \eqref{eq:4.42} only up to
\NSi{NS-F-P043-013}{X\in[Y_0,Y_1]} gives
\NSFormula{NS-F-P043-014}{display}{none}
\[
\max_{0\le j\le 2}\ \sup_{[Y_0,Y_1]\times[-1,1]}
\left|\partial_\eta^j(\mathfrak m_f-\mathfrak m_N)\right|\le C/N.
\]
By the shear formulas and Lemma~\ref{lem:4.4}\textup{(ii)}, increasing
\NSi{NS-F-P043-015}{N} if necessary gives
\NSFormula{NS-F-P043-016}{display}{none}
\[
E_f\ge \frac12\min_{[Y_0,Y_1]\times[-1,1]}E_0,
\qquad
\lVert a_f-a_N,b_f-b_N,p_f-p_N\rVert_0
\le C_{\mathrm{rep}}/N.
\]
Here \NSi{NS-F-P043-017}{C_{\mathrm{rep}}} is independent of
\NSi{NS-F-P043-018}{N}. Increase \NSi{NS-F-P043-019}{N} so that the final
triples stay in the fixed compact neighborhood used to define
\NSi{NS-F-P043-020}{C_\Psi}. Consequently \eqref{eq:4.41} remains strict on
the repair interval: for
\NSi{NS-F-P043-021}{N\ge 4C_\Psi C_{\mathrm{rep}}/\mu_R},
\NSFormula{NS-F-P043-022}{display}{none}
\[
\Psi_i(a_f,b_f,p_f)\ge \mu_R/2-C_\Psi C_{\mathrm{rep}}/N
\ge \mu_R/4>0.
\]
No profile or cumulative integral before \NSi{NS-F-P043-023}{Y_0} is affected
by this correction.

The equations \eqref{eq:4.42} give
\NSi{NS-F-P043-024}{\mathfrak m_f(Y_1)=\mathfrak m_0(Y_1)}. All profile changes are supported in
\NSi{NS-F-P043-025}{(X_-,Y_1)}; thus
\NSFormula{NS-F-P043-026}{display}{4.43}
\begin{equation}
\label{eq:4.43}
\begin{aligned}
(U_f,E_f)&=(U_0,E_0) &&\text{for }X\le X_-\text{ and }X\ge Y_1,\\
\mathfrak m_f(X)&=\mathfrak m_0(X) &&\text{for }X\ge Y_1,\\
(\Pi_f,V_f,Q_{s,f},N_{s,f},p_f,\mathcal T_{0,f})
&=(\Pi_0^{\mathrm{prof}},V_0,Q_{s,0},N_{s,0},p_0,
\mathcal T_{0,0})
&&\text{for }X\ge Y_1.
\end{aligned}
\tag{4.43}
\end{equation}
The last two equalities follow from Lemma~\ref{lem:4.4}\textup{(i)}. Here
\NSi{NS-F-P043-027}{V_f\text{ and }V_0} denote the radial-velocity profiles
computed from the final and joined pairs. All lower bounds on
\NSi{NS-F-P043-028}{N} encountered above are finite and depend only on already
fixed data; one finite integer can therefore satisfy them simultaneously. We
henceforth write \NSi{NS-F-P043-029}{(U,E,\Pi)} for the final profiles and
\NSi{NS-F-P043-030}{\mathcal T_0} for their stress.

\textbf{Step 4: verify the six conclusions.} The final profiles satisfy the cone
inequalities, and their five cumulative integrals agree exactly with the joined
pair beyond the correction interval. We now use these two facts to verify the
regularity, support, and endpoint estimates of the theorem. For the regularity
in part~\textup{(i)}, the changes in Steps 2--3 have compact support above
\NSi{NS-F-P043-031}{X_{\mathrm{an}}}. The analytic axis rectangle from
Proposition~\ref{prop:4.10} is therefore preserved. All remaining profile
pieces and their joins are smooth. Positivity of \NSi{NS-F-P043-032}{E}
follows from positivity of \NSi{NS-F-P043-033}{E_0}, the exponential in
\eqref{eq:4.38}, and the bound for \NSi{NS-F-P043-034}{E_f} on the repair
interval. Near the axis
\NSi{NS-F-P043-035}{E=\sqrt{2X}\phi/\mathsf C}; there
\NSi{NS-F-P043-036}{F=\phi/\mathsf C, U, \Pi, V_0/X} have the regularity asserted in
part~\textup{(i)}. On any finite rectangle this gives finite bounds for every
fixed mixed derivative.

The preserved pressure increment satisfies
\NSi{NS-F-P043-037}{C_p(\infty,\eta)=-\Pi_0(\eta)}. Hence
\NSFormula{NS-F-P043-038}{display}{none}
\[
\Pi(X,\eta)=\Pi_0(\eta)+C_p(X,\eta)
=C_p(X,\eta)-C_p(\infty,\eta)
=-\int_X^\infty\frac{E(x,\eta)^2}{2x}\,dx.
\]
This proves the pressure normalization in part~\textup{(i)} and the radial
balance in part~\textup{(ii)}. In particular
\NSi{NS-F-P043-039}{\Pi_X=E^2/(2X)=F^2} extends smoothly to the axis.
Proposition~\ref{prop:4.2} gives the exact
\NSPage{044}%
tangential residual identities. The first edge is unchanged, so the stress is
zero for \NSi{NS-F-P044-001}{X\le X_a} and \eqref{eq:4.13} holds there. By
\eqref{eq:4.43}, the outer stress is unchanged and zero for
\NSi{NS-F-P044-002}{X\ge X_b}.

On \NSi{NS-F-P044-003}{X_a<X<X_b}, the cone holds: it was preserved near the
two edges and beyond \NSi{NS-F-P044-004}{Y_1}, and was proved on
\NSi{NS-F-P044-005}{[X_-,Y_1]} in Steps 2--3. In particular, \eqref{eq:4.23}
gives
\NSFormula{NS-F-P044-006}{display}{none}
\[
\mathcal T_{0,\theta}+t_s\mathcal T_{0,z}=F(P_c-v_s)>0.
\]
Thus the stress is nonzero at every interior point, completing part~\textup{(ii)}.

We next prove the uniform assertions in part~\textup{(iii)}. Set
\NSi{NS-F-P044-007}{y_a=\log(X/X_a),\ y_b=\log(X_b/X)}. The factors
\eqref{eq:4.33}, \eqref{eq:4.32} give the following extensions of
\NSi{NS-F-P044-008}{n=\mathcal T_0/|\mathcal T_0|}:
\NSFormula{NS-F-P044-009}{display}{none}
\[
n=\frac{B_a}{|B_a|}\quad\text{near }X_a,
\qquad
n=\frac{(b_\theta,y_b^6b_z)}{\sqrt{b_\theta^2+y_b^{12}b_z^2}}
\quad\text{near }X_b.
\]
Their denominators stay positive on smaller collars. Since
\NSi{NS-F-P044-010}{B_a(0,\eta)=F(X_a,\eta)\mathfrak s(X_a,\eta)},
\NSFormula{NS-F-P044-011}{display}{none}
\[
n(X_a,\eta)=\frac{(1,t_s)}{\sqrt{1+t_s^2}},
\qquad n(X_b,\eta)=(1,0),
\qquad t_s(X_b,\eta)=0.
\]
The smooth functions \NSi{NS-F-P044-012}{F,a,v_s-2} are positive in the
interior; the endpoint estimates in Lemma~\ref{lem:4.9} and
Proposition~\ref{prop:4.10} make them positive also at
\NSi{NS-F-P044-013}{X_a,X_b}. They therefore have positive minima on the closed
annulus.

Define, on that rectangle,
\NSFormula{NS-F-P044-014}{display}{none}
\[
A_n=n_\theta+t_sn_z,
\qquad B_n=n_z-t_sn_\theta.
\]
In the interior, \eqref{eq:4.23} implies
\NSFormula{NS-F-P044-015}{display}{none}
\[
A_n=\frac{P_c-v_s}{|p_s-\mathfrak s|}>0.
\]
At the two edges, \NSi{NS-F-P044-016}{B_n=0}, while
\NSi{NS-F-P044-017}{A_n=\sqrt{1+t_s^2}} at \NSi{NS-F-P044-018}{X_a} and
\NSi{NS-F-P044-019}{A_n=1} at \NSi{NS-F-P044-020}{X_b}. Thus
\NSi{NS-F-P044-021}{A_n>0} on the closed rectangle, and we may define
\NSFormula{NS-F-P044-022}{display}{none}
\[
G_n=2-(v_s-2)B_n^2/A_n^2.
\]
By \eqref{eq:4.22},
\NSFormula{NS-F-P044-023}{display}{none}
\[
G_n=2-(v_s-2)J_c^2/(P_c-v_s)^2>0
\]
in the interior, while \NSi{NS-F-P044-024}{G_n=2} at both edges. Hence
\NSi{NS-F-P044-025}{A_n,G_n} are smooth and positive throughout. Taking
\NSFormula{NS-F-P044-026}{display}{none}
\[
\kappa=\min\left\{1,
\min_{[X_a,X_b]\times[-1,1]}A_n,
\min_{[X_a,X_b]\times[-1,1]}G_n\right\}>0
\]
gives \eqref{eq:4.26}, with \NSi{NS-F-P044-027}{\kappa<2}. This proves
part~\textup{(iii)} including its endpoint directions.

For part~\textup{(iv)}, let \NSi{NS-F-P044-028}{c_a=t_1^2>0} be the inner
exponential constant from \eqref{eq:4.33}, and put
\NSFormula{NS-F-P044-029}{display}{none}
\[
\zeta(X)=\exp(-c_a/y_a^2-4/y_b^2)\quad(X_a<X<X_b),
\qquad \zeta=0\quad\text{otherwise}.
\]
This function is smooth and flat at both edges. On a fixed inner collar, the
ratio \NSi{NS-F-P044-030}{\zeta/e^{-c_a/y_a^2}} is bounded above and below by
positive constants; on a fixed outer collar the same holds for
\NSi{NS-F-P044-031}{\zeta/e^{-4/y_b^2}}. The smooth nonzero factors in
\eqref{eq:4.33}, \eqref{eq:4.32} therefore give
\NSFormula{NS-F-P044-032}{display}{none}
\begin{align*}
|\mathcal T_0|&\ge c\zeta,& |\partial^\alpha\mathcal T_0|&\le C_\alpha\zeta y_a^{-m_\alpha}
&&\text{on the inner collar},\\
|\mathcal T_0|&\ge c\zeta,& |\partial^\alpha\mathcal T_0|&\le C_\alpha\zeta y_b^{-m_\alpha}
&&\text{on the outer collar}.
\end{align*}
Each derivative of an exponential factor contributes only finitely many inverse
powers of \NSi{NS-F-P044-033}{y_a\text{ or }y_b}. On the compact interval
between the collars, \NSi{NS-F-P044-034}{\zeta>0},
\NSi{NS-F-P044-035}{|\mathcal T_0|>0}, and all fixed derivatives are bounded. With
\NSi{NS-F-P044-036}{\delta=\min\{1,y_a,y_b\}}, these estimates combine to give
\eqref{eq:4.27}, proving part~\textup{(iv)}.
\NSPage{045}%
Finally, the total moment identities in part~\textup{(v)} survive both exact
connections. For the last correction this follows explicitly from
\NSFormula{NS-F-P045-001}{display}{none}
\begin{align*}
(M,J,S)_f(\infty)&=(M,J,S)_0(\infty)=(0,0,0),\\
\int_0^\infty(H_f-H_0)\,dX&=I_f(Y_1)-I_0(Y_1)=0.
\end{align*}
Together with Step 1, these equalities give \eqref{eq:4.28}. The exterior pair
remains the heat profile supplied by Lemma~\ref{lem:4.8}, with the same
\NSi{NS-F-P045-002}{c_\infty}. The identity
\NSFormula{NS-F-P045-003}{display}{none}
\[
q^{-A}c_\infty X^{-A}\mathcal H(2d/X)
=c_\infty s^{-A}\mathcal H(2\tau/s)
\]
proves the physical exterior formula. Since
\NSi{NS-F-P045-004}{Y_1<\sup\mathcal I_4<X_v}, equation \eqref{eq:4.43}
preserves \NSi{NS-F-P045-005}{U=M=J=0} for
\NSi{NS-F-P045-006}{X\ge X_v} with the same \NSi{NS-F-P045-007}{X_v} as in
Lemma~\ref{lem:4.8}. Equation \eqref{eq:4.7} yields
\NSFormula{NS-F-P045-008}{display}{none}
\[
V_0=\frac{2\eta XU-2D\eta M-dM_\eta}{L}=0.
\]
The radius \NSi{NS-F-P045-009}{X_v} precedes the terminal collar and lies after
all four reserved intervals. This proves part~\textup{(v)}. For
part~\textup{(vi)}, the heat compensation used \NSi{NS-F-P045-010}{\mathcal I_2},
and Step 3 used only \NSi{NS-F-P045-011}{\mathcal I_1}. Thus
\NSi{NS-F-P045-012}{\mathcal I_{\mathrm{pos}}=\mathcal I_3} and
\NSi{NS-F-P045-013}{\mathcal I_{\mathrm{mean}}=\mathcal I_4} retain the exact
profile \eqref{eq:4.30}, proving part~\textup{(vi)}. All choices were made in
profile coordinates before physical \NSi{NS-F-P045-014}{q} or any later
correction stage is introduced, as required.
\end{proof}

The profiles now give the leading field from Section~\ref{sec:3} with its
prescribed annular stress. The next section corrects the higher-order residual
terms identified after \eqref{eq:4.7}.

\section{Correcting the base flow to every order}
\label{sec:5}

Near the axis, the leading profiles
\NSi{NS-F-P045-015}{(E,U,V_0,\Pi)} of Theorem~\ref{thm:4.6} satisfy the balances
in Equations~\eqref{eq:4.7} and \eqref{eq:4.13}, but axial viscosity and the
remaining radial terms still contribute to the residual. We add axisymmetric
corrections to obtain the background \NSi{NS-F-P045-016}{(u_B,p_B)} described
in Section~\ref{sec:3} and Proposition~\ref{prop:5.5}. The velocity
\NSi{NS-F-P045-017}{u_B} is exactly divergence-free, and its momentum residual
is minus the cylindrical divergence of a physical annular stress
\NSi{NS-F-P045-018}{\mathcal T_{\mathrm{phys}}}, as in \eqref{eq:5.41}, plus an error
\NSi{NS-F-P045-019}{\mathcal E_B} bounded by every power of
\NSi{NS-F-P045-020}{q} as \NSi{NS-F-P045-021}{q\downarrow0}. The same bounds
hold for every fixed physical derivative on each compact profile range. The
annular stress is the input to the wave construction in Section~\ref{sec:7}.

These remaining terms enter at successive powers of
\NSi{NS-F-P045-022}{q^{2h}} in Equations~\eqref{eq:5.3} to \eqref{eq:5.6}. At
each order, Lemma~\ref{lem:5.1} solves for a correction near the axis, and
Lemma~\ref{lem:5.2} extends it radially and imposes the five integral conditions
that keep its stress supported in the prescribed annulus. The induction in
Proposition~\ref{prop:5.3} completes the current order before using its
coefficients to construct the next order.

The leading profiles and their construction constants remain fixed throughout.
The resulting coefficient sequence is the formal expansion \eqref{eq:5.1}:
convergence of the unmodified infinite series is not asserted. After
constructing this sequence, Lemma~\ref{lem:5.4} produces smooth fields for
\NSi{NS-F-P045-023}{q>0} with the same asymptotic coefficients. The proof of
Proposition~\ref{prop:5.5} checks that summation preserves incompressibility,
the exterior velocity, and the required residual estimates.

\subsection{Higher coefficients near the axis.}
\label{sec:5.1}
We first solve the coefficient equations Equations~\eqref{eq:5.2} to
\eqref{eq:5.6} near the axis, where the stress must remain zero. For
\NSi{NS-F-P045-024}{n\ge0}, we set
\NSi{NS-F-P045-025}{\lambda_n=2nh}, \NSi{NS-F-P045-026}{E_0=E},
\NSi{NS-F-P045-027}{U_0=U}, and \NSi{NS-F-P045-028}{\Pi_0=\Pi}. Each
coefficient profile is a function of \NSi{NS-F-P045-029}{(X,\eta)}. The
subscript zero now denotes expansion order:
\NSi{NS-F-P045-030}{\Pi_0(X,\eta)} is the entire leading pressure profile, whose
trace
\NSPage{046}%
\NSi{NS-F-P046-001}{\Pi_0(0,\eta)} is the axis pressure function defined in
\eqref{eq:4.31}. We use the similarity variables of \eqref{eq:4.1} and write
the physical coefficients as
\NSFormula{NS-F-P046-002}{display}{5.1}
\begin{equation}
\label{eq:5.1}
\begin{aligned}
u_{\theta,n}&=q^{-A+\lambda_n}E_n
=rq^{-A-1/2+\lambda_n}\phi_n/\mathsf C,
& E_n&=\sqrt{2X}\phi_n/\mathsf C,\\
u_{z,n}&=q^{-A+\lambda_n}U_n,
& ru_{r,n}&=q^{\lambda_n}V_n,
& p_n&=q^{-2A+\lambda_n}\Pi_n.
\end{aligned}
\tag{5.1}
\end{equation}
Here \NSi{NS-F-P046-003}{\phi_n}, rather than
\NSi{NS-F-P046-004}{E_n}, is the azimuthal coefficient that is smooth at
\NSi{NS-F-P046-005}{X=0}. The radial average is the operator from
\eqref{eq:4.6}, applied at fixed \NSi{NS-F-P046-006}{\eta}:
\NSi{NS-F-P046-007}{\mathcal A_X(U_n)(X,\eta)
=X^{-1}\int_0^XU_n(x,\eta)\,dx}.

The spacing \NSi{NS-F-P046-008}{2h} in \eqref{eq:5.1} accounts for the two
lower-order effects identified after \eqref{eq:4.7}. We recall
\NSi{NS-F-P046-009}{A=\tfrac12+h\text{ and }D=\tfrac12-h} from
\eqref{eq:4.1}, so \NSi{NS-F-P046-010}{1-2D=2A-1=2h}. First, the axial
derivative identity in \eqref{eq:4.2} gives the factor
\NSi{NS-F-P046-011}{q^{-2D}=q^{-1}q^{2h}} for two axial derivatives. The
radial relation \NSi{NS-F-P046-012}{X=r^2/(2q)} in \eqref{eq:4.1} gives
\NSi{NS-F-P046-013}{q^{-1}} for two radial derivatives. Thus axial viscosity
applied to order \NSi{NS-F-P046-014}{n-1} enters the angular and axial
equations at order \NSi{NS-F-P046-015}{n}; these are the
\NSi{NS-F-P046-016}{\text{order-}n-1} viscous terms in Equations~\eqref{eq:5.3}
and \eqref{eq:5.4} below.

Second, we recall the leading radial balance
\NSi{NS-F-P046-017}{r\partial_rp^{(0)}=(u_\theta^{(0)})^2} underlying the
pressure identity in \eqref{eq:4.7}. With the physical powers in
\eqref{eq:4.3}, both sides have power \NSi{NS-F-P046-018}{q^{-2A}}. The other
terms in the radial momentum equation multiplied by \NSi{NS-F-P046-019}{r}
start at \NSi{NS-F-P046-020}{q^{-1}=q^{-2A}q^{2h}}, as noted after
\eqref{eq:4.7}. This gives the order shift in the pressure equation
\eqref{eq:5.5}: the remaining radial terms are collected in \eqref{eq:5.6} and
enter one order later.

To include differentiation of the powers of \NSi{NS-F-P046-021}{q}, we recall
the operators \NSi{NS-F-P046-022}{T_a,Z_a} from \eqref{eq:4.2} and abbreviate
\NSFormula{NS-F-P046-023}{display}{none}
\[
T_{a,n}=T_{a+\lambda_n},
\qquad Z_{a,n}=Z_{a+\lambda_n},
\qquad Z_{a,n}^{[2]}=Z_{a+\lambda_n-D}Z_{a+\lambda_n}.
\]
The order of composition in the last expression accounts for the change of
power after the first axial derivative. Negative field indices mean zero.
Incompressibility, with the integration constant fixed by regularity at the
axis, is equivalent to
\NSFormula{NS-F-P046-024}{display}{5.2}
\begin{equation}
\label{eq:5.2}
\partial_XV_n=-Z_{-A,n}U_n,
\qquad
\frac{V_n}{X}
=\frac{2\eta U_n-2\eta(D+\lambda_n)\mathcal A_X(U_n)
-d\partial_\eta\mathcal A_X(U_n)}{L}.
\tag{5.2}
\end{equation}
We fix \NSi{NS-F-P046-025}{n\ge1}. At this point all coefficients of orders
\NSi{NS-F-P046-026}{j<n} are known, including their radial cutoffs and the
integral corrections constructed in Lemma~\ref{lem:5.2}. We solve for
\NSi{NS-F-P046-027}{\phi_n,U_n,\Pi_n} on a common inner interval;
\NSi{NS-F-P046-028}{V_n} is determined from \NSi{NS-F-P046-029}{U_n} by
\eqref{eq:5.2}. The pressure coefficient \NSi{NS-F-P046-030}{\Pi_n} must be
solved together with \NSi{NS-F-P046-031}{\phi_n,U_n}.

We set \NSi{NS-F-P046-032}{b=-A-\tfrac12} and
\NSi{NS-F-P046-033}{c=-A}. Primes in the following system mean
\NSi{NS-F-P046-034}{\partial_X}, with \NSi{NS-F-P046-035}{\eta} fixed. The
equations that cancel the order-\NSi{NS-F-P046-036}{n} momentum residual on
the inner interval are
\NSFormula{NS-F-P046-037}{display}{5.3}
\begin{equation}
\label{eq:5.3}
2(X\phi_n''+2\phi_n')
=T_{b,n}\phi_n
+\sum_{i+j=n}\left\{V_i(\phi_j'+\phi_j/X)+U_iZ_{b,j}\phi_j\right\}
-Z_{b,n-1}^{[2]}\phi_{n-1},
\tag{5.3}
\end{equation}
\NSFormula{NS-F-P046-038}{display}{5.4}
\begin{equation}
\label{eq:5.4}
2(XU_n''+U_n')
=T_{c,n}U_n
+\sum_{i+j=n}\left\{V_iU_j'+U_iZ_{c,j}U_j\right\}
+Z_{-2A,n}\Pi_n-Z_{c,n-1}^{[2]}U_{n-1},
\tag{5.4}
\end{equation}
\NSFormula{NS-F-P046-039}{display}{5.5}
\begin{equation}
\label{eq:5.5}
\Pi_n'=\mathsf C^{-2}\sum_{i+j=n}\phi_i\phi_j-\frac{\Omega_{n-1}}{2X},
\tag{5.5}
\end{equation}
\NSFormula{NS-F-P046-040}{display}{5.6}
\begin{equation}
\label{eq:5.6}
\Omega_k=T_{0,k}V_k
+\sum_{i+j=k}\left\{V_i(V_j'-V_j/(2X))+U_iZ_{0,j}V_j\right\}
-2XV_k''-Z_{0,k-1}^{[2]}V_{k-1}.
\tag{5.6}
\end{equation}
\NSPage{047}%
Every occurrence of \NSi{NS-F-P047-001}{\Omega_{n-1},\phi_{n-1},\text{ or }U_{n-1}}
in these source terms uses already constructed coefficients. For example, the
pressure equation separates into
\NSFormula{NS-F-P047-002}{display}{none}
\[
\Pi_n'=2\mathsf C^{-2}\phi_0\phi_n
+\mathsf C^{-2}\sum_{i=1}^{n-1}\phi_i\phi_{n-i}
-\frac{\Omega_{n-1}}{2X}.
\]
The sum is empty when \NSi{NS-F-P047-003}{n=1}. Its terms and
\NSi{NS-F-P047-004}{\Omega_{n-1}} are known, while the first term is linear in
the current unknown. Splitting the transport sums into
\NSi{NS-F-P047-005}{(i,j)=(0,n),(n,0)} and
\NSi{NS-F-P047-006}{1\le i,j<n} gives the same conclusion for the other
equations. Thus the system is linear at each positive order. The pressure
equation will also be imposed when the current profiles are extended beyond
the inner interval.

The powers of \NSi{NS-F-P047-007}{q} and the terms arising from the cylindrical
basis are obtained as follows. On
\NSi{NS-F-P047-008}{rq^{b+\lambda_n}\phi_n/\mathsf C}, angular radial transport
includes \NSi{NS-F-P047-009}{u_r(\partial_r+r^{-1})u_\theta} and therefore
yields \NSi{NS-F-P047-010}{V_i(\partial_X\phi_j+\phi_j/X)}. The angular
vector Laplacian gives
\NSi{NS-F-P047-011}{2(X\partial_{XX}\phi_n+2\partial_X\phi_n)} after removing
\NSi{NS-F-P047-012}{rq^{b-1+\lambda_n}/\mathsf C}. The corresponding axial scalar
Laplacian gives \NSi{NS-F-P047-013}{2(X\partial_{XX}U_n+\partial_XU_n)}.
Every transport product has the same power because
\NSi{NS-F-P047-014}{A+D=1}, and axial viscosity increases the order in the
\NSi{NS-F-P047-015}{q^{2h}} expansion by one because
\NSi{NS-F-P047-016}{1-2D=2h}. In the radial equation, multiplication by
\NSi{NS-F-P047-017}{r} places pressure and centrifugal acceleration at power
\NSi{NS-F-P047-018}{q^{-2A+\lambda_n}}, whereas the other radial terms start
at \NSi{NS-F-P047-019}{q^{-1+\lambda_k}}; hence
\NSi{NS-F-P047-020}{k=n-1}. Finally, applying the radial operator
\NSi{NS-F-P047-021}{\partial_{rr}+r^{-1}\partial_r-r^{-2}} to
\NSi{NS-F-P047-022}{V(X)/r} gives
\NSi{NS-F-P047-023}{2X\partial_{XX}V/(qr)}, giving the last two terms of
\eqref{eq:5.6}. Since \eqref{eq:5.2} gives
\NSi{NS-F-P047-024}{V_j=Xv_j} with smooth \NSi{NS-F-P047-025}{v_j}, every
term of \NSi{NS-F-P047-026}{\Omega_k} is divisible by \NSi{NS-F-P047-027}{X}.
For example,
\NSi{NS-F-P047-028}{V_i(\partial_XV_j-V_j/(2X))
=Xv_i(v_j/2+X\partial_Xv_j)}. Thus \NSi{NS-F-P047-029}{\Omega_k/X} is
regular at the axis, including all fixed parameter derivatives.

The linear equations must be solvable on one radial interval at every order.
This requires control of the \NSi{NS-F-P047-030}{\eta}-derivatives in the
recursion even as the coefficient bounds grow. The analytic dependence of the
leading inner profile supplies that control.

\begin{lemma}
\label{lem:5.1}
There exists an interval \NSi{NS-F-P047-031}{0\le\xi\le a},
\NSi{NS-F-P047-032}{\xi=\sqrt X}, extending from the axis into the inner collar
of Theorem~\ref{thm:4.6}\textup{(i)}, where the stress cone inequalities hold
with a uniform positive margin, on which every positive order has a unique
solution of Equations~\eqref{eq:5.2} to \eqref{eq:5.6} with
\NSi{NS-F-P047-033}{\phi_n(0,\eta)=U_n(0,\eta)=\Pi_n(0,\eta)=0}. The profiles
are smooth in \NSi{NS-F-P047-034}{X} and, for each fixed radial derivative,
holomorphic on a neighborhood of \NSi{NS-F-P047-035}{[-1,1]} in the parameter
\NSi{NS-F-P047-036}{\eta}. That neighborhood and its bounds may depend on the
order and the radial derivative; \NSi{NS-F-P047-037}{a} does not. Uniqueness
holds among solutions whose profiles extend holomorphically in
\NSi{NS-F-P047-038}{\eta} to a neighborhood of \NSi{NS-F-P047-039}{[-1,1]},
with bounds uniform for \NSi{NS-F-P047-040}{0\le\xi\le a}.
\end{lemma}

The zero axis values leave the leading traces of
\NSi{NS-F-P047-041}{\phi,U,\text{ and }\Pi} unchanged by every positive-order
coefficient.

\begin{proof}
\textbf{Step 1: write a linear system on the fixed inner interval.} The analytic
axis solution and its first collar, supplied by Proposition~\ref{prop:B.2} and
Theorem~\ref{thm:4.6}, have the stated parameter regularity on a fixed radial
rectangle. We choose \NSi{NS-F-P047-042}{a} inside this rectangle, beyond the
shorter collar that will remain uncut. At order \NSi{NS-F-P047-043}{n}, we use
the previously constructed lower-order coefficients, including their cutoffs
and moment corrections. The corrections are supported to the right of this
rectangle, and the radial cutoffs preserve parameter analyticity there. A
sufficiently small complex neighborhood
\NSi{NS-F-P047-044}{S_\rho=\{\eta\in\mathbb C:
\operatorname{dist}(\eta,[-1,1])<\rho\}} therefore contains no zero of
\NSi{NS-F-P047-045}{L} and supplies bounded holomorphic coefficients and
sources throughout \NSi{NS-F-P047-046}{0\le\xi\le a}.

We introduce the auxiliary variable \NSi{NS-F-P047-047}{K_n} and the vector
\NSi{NS-F-P047-048}{W_n} to write the system as
\NSFormula{NS-F-P047-049}{display}{5.7}
\begin{equation}
\label{eq:5.7}
\begin{aligned}
K_n&=\mathcal A_X(U_n)-U_n,
&W_n&=(\phi_n,U_n,K_n,\Pi_n,\partial_\xi\phi_n,\partial_\xi U_n)^{\mathsf T},\\
\partial_\xi W_n+\xi^{-1}\operatorname{diag}(0,0,2,0,3,1)W_n
&=A_0W_n+A_1\partial_\eta W_n+f_n,
&W_n(0)&=0.
\end{aligned}
\tag{5.7}
\end{equation}
\NSPage{048}%
The matrices \NSi{NS-F-P048-001}{A_0,A_1} are
\NSi{NS-F-P048-002}{6\times6} coefficient matrices in
\NSi{NS-F-P048-003}{(\xi,\eta)}, with \NSi{NS-F-P048-004}{n} fixed;
\NSi{NS-F-P048-005}{f_n} is a known six-component source formed from the
lower-order profiles. The additional unknown \NSi{NS-F-P048-006}{K_n}
replaces the radial average by a first-order equation, so the system has no
unevaluated radial integral of the current unknowns. Indeed,
\NSi{NS-F-P048-007}{\partial_\xi K_n+2K_n/\xi=-\partial_\xi U_n}, and the two
radial second-order operators become
\NSi{NS-F-P048-008}{(\partial_{\xi\xi}+3\xi^{-1}\partial_\xi)/2} and
\NSi{NS-F-P048-009}{(\partial_{\xi\xi}+\xi^{-1}\partial_\xi)/2}. At positive
order all terms containing an order-\NSi{NS-F-P048-010}{n} unknown are linear.
The pressure row is
\NSi{NS-F-P048-011}{\partial_\xi\Pi_n=4\xi\phi_0\phi_n/\mathsf C^2+
\text{known terms}}; we also use this row for
\NSi{NS-F-P048-012}{X\partial_X\Pi_n} in the axial equation. All apparent
divisions by \NSi{NS-F-P048-013}{X} are then regular.

\textbf{Step 2: solve without shrinking the radial interval.} Iteration of the
system differentiates in \NSi{NS-F-P048-014}{\eta}. We control the number of
such derivatives before estimating the successive radial integrals. The
coefficient of \NSi{NS-F-P048-015}{\partial_\eta} has the following sparse
block form. With \NSi{NS-F-P048-016}{H_c=D\eta+dU_0}, its only possibly
nonzero entries are
\NSFormula{NS-F-P048-017}{display}{none}
\begin{align*}
(A_1)_{51}&=2H_c/L,
& (A_1)_{52}=(A_1)_{53}&=-2d(\phi_0+X\partial_X\phi_0)/L,\\
(A_1)_{62}&=2(H_c-dX\partial_XU_0)/L,
& (A_1)_{63}&=-2dX\partial_XU_0/L,
& (A_1)_{64}&=2d/L.
\end{align*}
Thus \NSi{NS-F-P048-018}{A_1}, at every radius and after every parameter
derivative, maps the first four coordinates into the last two and annihilates
the last two.

For \NSi{NS-F-P048-019}{c_i=(0,0,2,0,3,1)_i}, we invert the diagonal part of
\eqref{eq:5.7}, with zero axis data, by the integral operator
\NSFormula{NS-F-P048-020}{display}{none}
\[
(Gg)_i(\xi)=\int_0^\xi(s/\xi)^{c_i}g_i(s)\,ds,
\qquad \mathcal K=G(A_0+A_1\partial_\eta).
\]
Its kernels are independent of \NSi{NS-F-P048-021}{\eta} and preserve the two
coordinate subspaces just described. If \NSi{NS-F-P048-022}{D_0} is any
intervening diagonal kernel, then
\NSFormula{NS-F-P048-023}{display}{none}
\[
A_1(\xi)D_0A_1(s)=0,
\qquad A_1(\xi)D_0\partial_\eta A_1(s)=0.
\]
Consequently two adjacent derivative operators vanish even when the left
derivative falls on a variable matrix coefficient. A nonzero composition of
\NSi{NS-F-P048-024}{k} factors contains at most
\NSi{NS-F-P048-025}{p_k=\lceil k/2\rceil} parameter derivatives.

We bound \NSi{NS-F-P048-026}{A_0,A_1,f_n} on the larger neighborhood
\NSi{NS-F-P048-027}{S_\rho}, uniformly on the whole radial interval. On
\NSi{NS-F-P048-028}{S_{\rho'}}, \NSi{NS-F-P048-029}{0<\rho'<\rho}, we split
the Cauchy radius loss \NSi{NS-F-P048-030}{\Delta=\rho-\rho'} among the at
most \NSi{NS-F-P048-031}{p_k} derivatives. The diagonal kernels have norm at
most one, and the ordered integration simplex has volume
\NSi{NS-F-P048-032}{a^{k+1}/(k+1)!}. Summing the at most
\NSi{NS-F-P048-033}{2^k} compositions gives
\NSFormula{NS-F-P048-034}{display}{5.8}
\begin{equation}
\label{eq:5.8}
\lVert\mathcal K^kGf_n\rVert_{S_{\rho'}}
\le \frac{C_n^{k+1}a^{k+1}}{(k+1)!}
\left(\max\{1,p_k/\Delta\}\right)^{p_k}.
\tag{5.8}
\end{equation}
The \NSi{NS-F-P048-035}{k}th root of the right-hand side of \eqref{eq:5.8}
tends to zero. The Picard series
\NSi{NS-F-P048-036}{W_n=\sum_{k\ge0}\mathcal K^kGf_n} therefore converges on
\NSi{NS-F-P048-037}{S_{\rho'}} for every finite \NSi{NS-F-P048-038}{a},
regardless of the size of \NSi{NS-F-P048-039}{C_n}. Iterating the homogeneous
equation for the difference of two solutions and applying the same bound proves
uniqueness. In particular, the radial interval is independent of the
coefficient norms at later orders.

\textbf{Step 3: recover smooth profiles at the axis.} Radial smoothness follows
without differentiating a singular expression:
\NSFormula{NS-F-P048-040}{display}{none}
\[
(Gg)_i=\xi\int_0^1t^{c_i}g_i(t\xi)\,dt,
\qquad
\partial_\xi(Gg)_i=g_i(\xi)-c_i\int_0^1t^{c_i}g_i(t\xi)\,dt.
\]
The second identity is continuous at zero. Repeated use of it, with a smaller
parameter neighborhood when necessary, gives every fixed radial derivative.
We extend the integral
